% RESEARCH_DRAFT_SHA256: 1092844041354d4b31e4a5019166eb1a95c5304592dadf4e6bd6c791dbcab027
% Layout and source-reference wrappers expanded for the web renderer.
\documentclass[11pt]{article}
\usepackage[margin=25mm]{geometry}
\usepackage{fontspec,unicode-math}
\setmainfont{Latin Modern Roman}
\setmathfont{Latin Modern Math}
\usepackage{amsmath,amssymb,mathtools,xurl,hyperref}
\hypersetup{colorlinks=true,urlcolor=blue,linkcolor=blue}
\setcounter{secnumdepth}{0}
\setlength{\parindent}{0pt}
\setlength{\parskip}{5pt}
\setlength{\emergencystretch}{3em}
\newcommand{\sref}[2]{\hyperlink{source:#1:#2}{[S#2]}}
\newcommand{\cataloguescope}[1]{\paragraph{Source record.}#1}
\newcommand{\floor}[1]{\left\lfloor#1\right\rfloor}
\begin{document}
\section*{Research-attempt provenance}
This standalone research attempt retains the complete original definition body below,
including its historical source and preparation notes. Those notes describe the
earlier catalogue preparation; the new attempt and its verification are identified
separately after that body. No original ID, catalogue status, or source statement is changed.

\section{A joint tree and JC69/K80 selector with a finite-sample certificate}\label{statement}

\subsection{Definitions and mathematical statement}

\paragraph{Shared notation.}Standard number systems are $\mathbb N,\mathbb Z,\mathbb Q,\mathbb R,\mathbb C$. Specialized notation and index ranges follow the cited source and the definitions below.


A phylogenetic model assigns a probability distribution to observed leaf states on a tree using hidden internal states and stochastic transition matrices along edges. A phylogenetic invariant is a polynomial vanishing on that model. Flattening groups tensor indices into blocks to produce a matrix or lower-order tensor; rank constraints can then give polynomial minors. The exact constraints depend on the mutation model and state cardinality. Tree fitting compares a candidate tree's predicted distribution with observed data using a specified criterion, such as likelihood. Mutation-model selection compares different restrictions on edge transitions.
\paragraph{Record scope and qualifications.}
The catalogue title is retained as metadata. Named certificates, contracts, bounds or frameworks in that title are not asserted as results of the original source extract.
The excerpt specifies neither the candidate tree class nor an accuracy or fitting criterion. No JC69/K80 selection procedure is inferred from the expanded catalogue title. The supplied catalogue marks this record partially solved; its present research status was not independently checked.
The source text below is rendered literally where necessary; damaged formulas and composite extracts are not silently repaired into a new mathematical claim.
\paragraph{Statement recorded in the supplied source.}
Question 33. How can we efficiently decide which trees fit the data best? Which model of mutation is most accurate? \href{https://aimath.org/WWN/compalgstat/compalgstat.pdf}{[S2]}
\paragraph{Source status.}
The supplied catalogue labels this source record partially solved; that metadata was checked against the supplied snapshot on 2026-10-11. The label is preserved for provenance and is not an independent verification of present open status or mathematical correctness.
\subsection{Short English statement}
The source asks how to choose trees and mutation models efficiently.
\subsection{Sources}
\begin{itemize}
\item[S1] Ulam AI and the UnsolvedMath Contributors, UnsolvedMath: A Curated Collection of Open Mathematics Problems, catalogue entry. Catalogue curation and metadata: CC BY 4.0. The source extract is retained; definitions, English summaries and qualifications are editorial additions.
\url{https://huggingface.co/datasets/ulamai/UnsolvedMath}.
\item[S2] AIM Workshop Problem Lists, source linked by the catalogue.
\url{https://aimath.org/WWN/compalgstat/compalgstat.pdf}.
\end{itemize}
\paragraph{Preparation and review.}
The original source and UnsolvedMath contributors retain attribution for the source extract. Codex (GPT-6; exact serving revision not exposed) assisted with definitions, English summaries, scope notes and formatting. No new human mathematical authorship or independently verified solution is asserted. Catalogue identity checked against the complete supplied registry and current repository. Distinct upstream records, including repeated questions and broad or damaged extracts, are retained for complete coverage of the source CS category. A repeated record is not a new independent conjecture. Complete mathematical-equivalence and current-literature review has not been performed. Preparation and final review follow contribution prompts C and D, with source gaps disclosed. Structural and rendering checks do not establish a new theorem or open status.
\label{end:statement}

\section{Literature and context check}
The original AIM Computational Algebraic Statistics
 workshop PDF was consulted. The preserved extract does not specify every model,
 parameterization or performance criterion needed by its broad question. The
 following declared model is an actual test of one route, not a replacement
 target or a claim of novelty for elementary statistical or algebraic facts.
\section{Research attempt: Common model distributions prevent uniform tree selection}
Choose a finite-state general Markov model on any tree, and let every edge
transition matrix have all rows equal to the same strictly positive probability
vector $\pi$. Each child is then drawn from $\pi$ independently of its parent.
Multiplying the conditional probabilities and summing hidden variables shows
the leaf distribution is $\prod_{i\in\text{leaves}}\pi(x_i)$, independently
of the tree. This holds for each of the three distinct four-leaf binary
topologies, even though all transition entries are positive.

Given any selector based on $N$ independent sites, its output has the same law
under these three parameterized trees. The probabilities of outputting the
three labels sum to at most one, so at least one tree has success probability
at most $1/3$. No finite sample size, or even access to the exact leaf
distribution, identifies the originating topology at these parameters.
This is an identifiability obstruction, not a computational hardness result.

A quantitative version explains why nearby regimes need a separation margin.
For two distributions $p,q$ on a finite alphabet, define
$d_{TV}(p,q)=\frac12\sum_x|p(x)-q(x)|$. Expanding a product difference one
factor at a time and summing absolute values gives
$d_{TV}(p^{\otimes N},q^{\otimes N})\le N d_{TV}(p,q)$.
For any decision event $A$, the sum of its two classification errors is
$p^{\otimes N}(A^c)+q^{\otimes N}(A)
\ge1-d_{TV}(p^{\otimes N},q^{\otimes N})$.
The event bound follows because the sum of positive entries of $p-q$ equals
its total variation. Thus arbitrarily close competing distributions preclude
a uniform finite-sample guarantee at fixed $N$.

For mutation models nested by parameter restriction, data generated at their
intersection similarly do not identify an originating model label. A criterion
such as the smallest adequate model is a separate convention. The source's
``best'' and ``most accurate'' cannot be replaced by unique label recovery
without such choices.
\section{Verification and adversarial audit}
The leaf-independence construction was checked by summing normalized
hidden factors, and preserves positive transition entries. It does allow
singular transition matrices; restrictions to invertible or specific
continuous-time transitions would be additional assumptions. Product-law and
testing inequalities were proved directly.
\paragraph{Audit corrections.} Positive transition entries were not confused with identifiability; the exact unrestricted-model obstruction was not imposed on every specialized mutation model.
\paragraph{External computation.} None; symbolic arguments and boundary checks only.
\section{Scope and remaining gaps}
The candidate class, mutation restrictions and fitting criterion are absent.
The first gap is a specified identifiable comparison with quantitative separation
and a proved algorithmic guarantee. No universal tree or mutation selector is
provided by the obstruction.
\paragraph{Integration review.} Codex (GPT-6; exact serving revision not exposed) reviewed the written arguments and expanded layout and source-reference wrappers. Repository metadata and the production web renderer were checked using Python 3.13 and JavaScript with MathJax 3.2.2. This model-assisted review is not independent human proof verification; no PDF compilation or formal proof checking is asserted.

\section{Completion Estimate}
\noindent\textbf{COMPLETION: 0\%}
The calculations below the preserved source statement investigate a concrete route,
but establish no substantiated new progress toward the entire original target.
Known special cases and failed sufficient conditions are not credited as a solution.
\section{References}
The original source references and catalogue attribution are preserved in the
Sources section above. The following URLs were actually consulted for this attempt;
retention of an additional original reference does not claim fresh consultation of it.
\begin{itemize}
\item \url{https://aimath.org/WWN/compalgstat/compalgstat.pdf} (consulted 2026-10-11). 
\end{itemize}
\section{Model and contributor disclosure}
OpenAI gpt-6-astra, requested reasoning configuration ultra, generated this research
attempt on 2026-10-11. The exact serving revision is not exposed. Codex GPT-6
performed orchestration, repository-format review, and a model-assisted mathematical
audit of the written arguments; its exact serving revision
is also not exposed. The mathematical route, boundary checks and final adversarial
self-audit followed contribution prompts A then B in
\texttt{docs/contribution\_prompts.md}, after reading \texttt{CONTRIBUTING.md}.
Python 3.14.2 assembled the standalone documents; any mathematical computation is
specified in Verification. No theorem prover was used. No human mathematical
authorship, substantial human proof correction, or independent human proof checking
is asserted. Model-assisted review, source retrieval, and structural checks do not
establish correctness or novelty. Full private reasoning transcripts are unavailable
as reproducibility artifacts; the complete arguments, inputs and outcomes actually
used are stated here or in the companion batch audit.
\section{Final status}
Unresolved as a new research contribution. This is a concrete route investigation,
not a claim of new complete resolution or a current-status certification. A known
answer reproduced for context does not become a new contribution. The catalogue's
source-status label is preserved; this attempt status does not assert that a
source-labelled solved question is open.
\end{document}
